\chapter{AUTONOMOUS MODE: SEARCH, MAP AND CHASE}
\label{ch:autonomy}

The goal of the project is a robot that, once switched on, finds the dog (or
the kids) on its own and chases her. Autonomous mode puts the pieces of the
previous chapters together: lidar mapping and exploration
(\S\ref{sec:explore}), localization in a saved map, and the chase controller
of Chapter~\ref{ch:chase}. It is the robot's default behaviour at boot.
The code is the \texttt{bbai\_autonomy} package
(\file{ros/bbai\_autonomy/}).

\section{BEHAVIOUR}
\label{sec:autobehave}

\begin{description}
\item[Explore] (first run, no saved map). gmapping builds the map and
  explore\_lite drives frontier to frontier, exactly as in
  \S\ref{sec:explore}. When explore\_lite has sent no new goal for
  \SI{60}{s} (and at least \SI{120}{s} have passed), the house is mapped.
  The mission saves it as \file{maps/house\_<date>.\{pgm,yaml\}}, points
  \file{maps/house.yaml} at it and restarts the host side in patrol mode.
  A partial map is also saved every \SI{2}{min} in case the power fails.
\item[Patrol] (every run after that). map\_server loads
  \file{house.yaml} and AMCL localizes the robot in it. The mission spreads
  search points over the map (\S\ref{sec:waypoints}), visits them nearest
  first, and at each one turns twice by \SI{120}{\degree} so the camera's
  \SI{39}{\degree} view sweeps the whole room. When every point has been
  visited it starts the next round, in the opposite order, forever.
\item[Chase] (any time). The chase node runs in \emph{auto} arm mode: it
  drives toward a target whenever it holds a track, without R2. While
  \topic{/chase/active} is true the mission cancels its patrol goal; once
  the target has been gone for \SI{3}{s} it re-sends the same goal. During
  exploration move\_base simply keeps planning in the background and
  continues after the chase.
\end{description}

\section{ARCHITECTURE}
\label{sec:autoarch}

The work is split by where it has to run. Anything that touches the motors or
the camera stays on the board; anything heavy runs on an off-board
\emph{host}, which is either the Proxmox server VM or the laptop. The host is
whatever machine sources \file{\textasciitilde/bbai\_slam/env.sh}; nothing in
the package names one.

\begin{table}[ht]
\centering\small
\caption{Autonomous-mode nodes.}
\label{tab:autonodes}
\begin{tabularx}{\linewidth}{@{}lXl@{}}
\toprule
Machine & Nodes & Started by \\
\midrule
Board (Melodic) & \texttt{robot.launch}, PS4 teleop, \texttt{chase} (auto),
  \texttt{cmd\_mux} & \texttt{bbai-robot.service} \\
Host (Noetic) & gmapping + explore\_lite \emph{or} map\_server + AMCL;
  move\_base; \texttt{mission} & \texttt{bbai-host@<user>.service} \\
\bottomrule
\end{tabularx}
\end{table}

\paragraph{Who drives.} Three sources want \topic{/cmd\_vel}. Rather than let
them fight (the problem noted for exploration in \S\ref{sec:explore}), each
publishes to its own topic and \texttt{cmd\_mux} on the board forwards one of
them, in strict priority:
\begin{enumerate}[nosep]
\item PS4 teleop while L1 is held (the mux goes silent; teleop publishes
  \topic{/cmd\_vel} itself);
\item chase, while \topic{/chase/cmd\_vel} is fresher than \SI{0.4}{s};
\item navigation, \topic{/nav/cmd\_vel} from move\_base, fresher than
  \SI{0.6}{s}.
\end{enumerate}
Nothing passes while autonomy is disabled. The Options button toggles it,
as does \topic{/autonomy/enable}. Each hand-over to nobody sends one zero
twist.

\paragraph{Failure behaviour.} If the host crashes or the Wi-Fi drops,
\topic{/nav/cmd\_vel} stops and the base node's \SI{0.5}{s} timeout halts the
motors. The board still teleops and still chases, because chase never depended
on the host. When the board reboots it starts a new roscore; the host script
watches the master's \texttt{run\_id} and restarts its side, so the two
re-pair without anyone logging in.

\section{TARGETS AND THE KID CLASS}
\label{sec:kids}

Targets are chosen at launch (\texttt{targets:=[dog,kid]}) or live with
\topic{/chase/set\_targets}. Any MobileNet-SSD label works
(\S\ref{sec:mobilenet}), plus \texttt{kid} and \texttt{adult}.

MobileNet-SSD has a \emph{person} class but no child class, so a person is
sorted by the height of the top of their head above the floor. With the box's
top edge at normalized image row $y_{\min}$ (0 at the top), depth $Z$ from the
spatial detector, RGB focal length $f_y$ at \SI{1080}{p}, camera height $h_c$
and upward pitch $\theta_c$,
\begin{equation}\label{eq:headheight}
h_{\mathrm{head}} = h_c + Z\,\tan\!\Big(\arctan\frac{(0.5-y_{\min})\,1080}{f_y}+\theta_c\Big).
\end{equation}
A person with $h_{\mathrm{head}}<\SI{1.35}{m}$ is a kid, otherwise an adult.
Only the top of the box is used, so the feet need not be in view. That matters
with the camera this low: with $h_c\approx\SI{0.15}{m}$ and a vertical field of
view near \SI{40}{\degree}, a \SI{1.1}{m} child fits in the frame head to toe
only beyond about \SI{2.6}{m}.

If the box touches the top of the frame the head is out of view, and
\eqref{eq:headheight} at $y_{\min}=0$ is only a lower bound. Above
\SI{1.35}{m} that still proves an adult; below it the node can't tell and does
not start a chase. Once it is already following a kid, such views keep the
track, since a kid up close always fills the frame. The status line prints the
estimate (\texttt{kid(1.12m tall)}) so $h_c$, $\theta_c$ and the threshold can
be tuned with real people in view. The low camera works in our favour here:
an adult's head is usually out of frame and above the threshold at the frame
top, while a child's head is in view from a couple of metres away.

\section{SAVED MAP OR A FRESH MAP EACH TIME}
\label{sec:savedmap}

The robot maps once and localizes after that. A fresh map each run was
rejected for three reasons:
\begin{itemize}[nosep]
\item AMCL on a fixed map is lighter and steadier than gmapping, and a long
  chase around one room cannot smear the map.
\item A finished map gives search points over every room. Frontier
  exploration stops as soon as no frontiers are left, which is exactly when a
  search needs to continue.
\item Mapping with the RPLIDAR A1 is slow and sometimes fails
  (Chapter~\ref{ch:slam}). Doing it once and checking the result beats
  redoing it on every boot.
\end{itemize}
The cost is that AMCL needs a starting guess. The mission writes the robot's
pose to \file{maps/last\_pose.yaml} every \SI{5}{s} and uses it at the next
start, so a robot switched off where it stopped resumes correctly. If it was
carried elsewhere, the command \texttt{home} declares it at the map origin
(where mapping began) and \texttt{global} spreads AMCL's particles over the
whole map. If the furniture moves, \texttt{remap} forgets the map and the
next run explores again.

AMCL uses the likelihood-field laser model with 90 of the 360 beams, 300 to
2000 particles (cheap on the host), an update every \SI{0.1}{m} or
\SI{0.2}{rad}, and a raised rotation noise ($\alpha_1=0.2$) for the skid
steer's slip when turning.

\section{SEARCH POINTS}
\label{sec:waypoints}

From the occupancy grid $M$ with resolution $r=\SI{0.05}{m}$:
\begin{enumerate}[nosep]
\item Mark cells free when $0\le M<25$; unknown and occupied cells are not.
\item Erode the free set by $k=\lceil 0.30/r\rceil$ cells (8-connected), so
  every remaining cell is at least \SI{0.30}{m} from a wall or unknown space.
\item Flood-fill from the cell nearest the robot. Only reachable cells count,
  which drops closed rooms and space seen through doorways that are too
  narrow.
\item Cut the map into $\SI{1.5}{m}\times\SI{1.5}{m}$ blocks and keep, in
  each block that has a reachable cell, the one nearest the block centre.
\item Order the points as a nearest-neighbour tour from the robot.
\end{enumerate}
A goal that fails or takes longer than \SI{90}{s} is skipped; the next round
tries it again.

\section{RUNNING IT}

\begin{description}
\item[By hand.] Board:
  \texttt{robot\_start.sh autostart:=false}. Host:
  \texttt{host\_autonomy.sh}. Press Options to start and again to stop.
\item[At boot.] Enable \texttt{bbai-robot.service} on the board and
  \texttt{bbai-host@<user>} on the host. Boot options go in
  \file{/etc/default/bbai-robot}, for example
  \texttt{BBAI\_ARGS="targets:=[dog,kid] autostart:=false"}; booting with
  autonomy off and starting it with Options is the safer default until it has
  run a few times.
\item[Watching.] \topic{/autonomy/state} says what the mission is doing
  (\texttt{PATROL round 2 point 5/14}) and \topic{/autonomy/source} who is
  driving.
\end{description}
The exact commands are in \file{docs/COMMANDS.md}.

\paragraph{Limits.} Everything listed for exploration in \S\ref{sec:explore}
still applies, stairs above all. The chase controller has no obstacle
avoidance of its own: it drives straight at the target through anything the
camera does not report. Its speed cap (\SI{0.2}{m/s}) and stop distance
(\SI{0.5}{m}) are the protection.
\influx{Autonomous mode is written and unit-tested off the robot (mux
priorities, kid classification, search points on a synthetic map). It has not
yet run on the robot.}
